%%
%% This is file `sample-sigconf-authordraft.tex',
%% generated with the docstrip utility.
%%
%% The original source files were:
%%
%% samples.dtx  (with options: `all,proceedings,bibtex,authordraft')
%% 
%% IMPORTANT NOTICE:
%% 
%% For the copyright see the source file.
%% 
%% Any modified versions of this file must be renamed
%% with new filenames distinct from sample-sigconf-authordraft.tex.
%% 
%% For distribution of the original source see the terms
%% for copying and modification in the file samples.dtx.
%% 
%% This generated file may be distributed as long as the
%% original source files, as listed above, are part of the
%% same distribution. (The sources need not necessarily be
%% in the same archive or directory.)
%%
%%
%% Commands for TeXCount
%TC:macro \cite [option:text,text]
%TC:macro \citep [option:text,text]
%TC:macro \citet [option:text,text]
%TC:envir table 0 1
%TC:envir table* 0 1
%TC:envir tabular [ignore] word
%TC:envir displaymath 0 word
%TC:envir math 0 word
%TC:envir comment 0 0
%%
%% The first command in your LaTeX source must be the \documentclass
%% command.
%%
%% For submission and review of your manuscript please change the
%% command to \documentclass[manuscript, screen, review]{acmart}.
%%
%% When submitting camera ready or to TAPS, please change the command
%% to \documentclass[sigconf]{acmart} or whichever template is required
%% for your publication.
%%
%%
%\documentclass[sigconf,authordraft]{acmart}
\documentclass[sigconf,anonymous,review]{acmart}
%%
%% \BibTeX command to typeset BibTeX logo in the docs
\AtBeginDocument{%
  \providecommand\BibTeX{{%
    Bib\TeX}}}

%% Rights management information.  This information is sent to you
%% when you complete the rights form.  These commands have SAMPLE
%% values in them; it is your responsibility as an author to replace
%% the commands and values with those provided to you when you
%% complete the rights form.
\setcopyright{acmlicensed}
\copyrightyear{2026}
\acmYear{2026}
\acmDOI{XXXXXXX.XXXXXXX}
%% These commands are for a PROCEEDINGS abstract or paper.
\acmConference[Conference acronym 'XX]{Make sure to enter the correct
  conference title from your rights confirmation email}{Month DD--DD,
  2026}{City, State, Country}
%%
%%  Uncomment \acmBooktitle if the title of the proceedings is different
%%  from ``Proceedings of ...''!
%%
%%\acmBooktitle{Woodstock '18: ACM Symposium on Neural Gaze Detection,
%%  June 03--05, 2018, Woodstock, NY}
\acmISBN{978-1-4503-XXXX-X/2026/XX}


%%
%% Submission ID.
%% Use this when submitting an article to a sponsored event. You'll
%% receive a unique submission ID from the organizers
%% of the event, and this ID should be used as the parameter to this command.
%%\acmSubmissionID{123-A56-BU3}

%%
%% For managing citations, it is recommended to use bibliography
%% files in BibTeX format.
%%
%% You can then either use BibTeX with the ACM-Reference-Format style,
%% or BibLaTeX with the acmnumeric or acmauthoryear sytles, that include
%% support for advanced citation of software artefact from the
%% biblatex-software package, also separately available on CTAN.
%%
%% Look at the sample-*-biblatex.tex files for templates showcasing
%% the biblatex styles.
%%

%%
%% The majority of ACM publications use numbered citations and
%% references.  The command \citestyle{authoryear} switches to the
%% "author year" style.
%%
%% If you are preparing content for an event
%% sponsored by ACM SIGGRAPH, you must use the "author year" style of
%% citations and references.
%% Uncommenting
%% the next command will enable that style.
%%\citestyle{acmauthoryear}


%%
%% end of the preamble, start of the body of the document source.
\begin{document}

%%
%% The "title" command has an optional parameter,
%% allowing the author to define a "short title" to be used in page headers.
\title{Multi-agent System for end-to-end Feature Delivery}

%%
%% The "author" command and its associated commands are used to define
%% the authors and their affiliations.
%% Of note is the shared affiliation of the first two authors, and the
%% "authornote" and "authornotemark" commands
%% used to denote shared contribution to the research.
\author{Ben Trovato}
\authornote{Both authors contributed equally to this research.}
\email{trovato@corporation.com}
\orcid{1234-5678-9012}
\author{G.K.M. Tobin}
\authornotemark[1]
\email{webmaster@marysville-ohio.com}
\affiliation{%
  \institution{Institute for Clarity in Documentation}
  \city{Dublin}
  \state{Ohio}
  \country{USA}
}

\author{Lars Th{\o}rv{\"a}ld}
\affiliation{%
  \institution{The Th{\o}rv{\"a}ld Group}
  \city{Hekla}
  \country{Iceland}}
\email{larst@affiliation.org}

\author{Valerie B\'eranger}
\affiliation{%
  \institution{Inria Paris-Rocquencourt}
  \city{Rocquencourt}
  \country{France}
}

\author{Aparna Patel}
\affiliation{%
 \institution{Rajiv Gandhi University}
 \city{Doimukh}
 \state{Arunachal Pradesh}
 \country{India}}

\author{Huifen Chan}
\affiliation{%
  \institution{Tsinghua University}
  \city{Haidian Qu}
  \state{Beijing Shi}
  \country{China}}

\author{Charles Palmer}
\affiliation{%
  \institution{Palmer Research Laboratories}
  \city{San Antonio}
  \state{Texas}
  \country{USA}}
\email{cpalmer@prl.com}

\author{John Smith}
\affiliation{%
  \institution{The Th{\o}rv{\"a}ld Group}
  \city{Hekla}
  \country{Iceland}}
\email{jsmith@affiliation.org}

\author{Julius P. Kumquat}
\affiliation{%
  \institution{The Kumquat Consortium}
  \city{New York}
  \country{USA}}
\email{jpkumquat@consortium.net}

%%
%% By default, the full list of authors will be used in the page
%% headers. Often, this list is too long, and will overlap
%% other information printed in the page headers. This command allows
%% the author to define a more concise list
%% of authors' names for this purpose.
\renewcommand{\shortauthors}{Trovato et al.}

%%
%% The abstract is a short summary of the work to be presented in the
%% article.
\begin{abstract}
Spec-driven development (SDD) with AI coding agents provides a structured workflow, but agents often remain ``context blind'' in large, evolving repositories, leading to hallucinated APIs and violations of local architecture. We present \textbf{Spec Kit Agents}, a compound multi-agent system built on the Spec Kit workflow that adds a \textbf{tool-augmented guardrail layer}. The layer uses read-only prober agents to ground each stage of development (Specify, Plan, Task, Implement) in repository evidence and to validate plans against the local environment. In our evaluation, these guardrails improve correctness while reducing delivery time for complex features by up to 48.7\% by pruning invalid implementation paths early.
\end{abstract}

%%
%% The code below is generated by the tool at http://dl.acm.org/ccs.cfm.
%% Please copy and paste the code instead of the example below.
%%
\begin{CCSXML}
<ccs2012>
 <concept>
  <concept_id>10011007.10011006.10011066.10011069</concept_id>
  <concept_desc>Software and its engineering~Software development techniques</concept_desc>
  <concept_significance>500</concept_significance>
 </concept>
 <concept>
  <concept_id>10011007.10011006.10011039.10011040</concept_id>
  <concept_desc>Software and its engineering~Software verification and validation</concept_desc>
  <concept_significance>300</concept_significance>
 </concept>
 <concept>
  <concept_id>10010147.10010178.10010219.10010220</concept_id>
  <concept_desc>Computing methodologies~Multi-agent systems</concept_desc>
  <concept_significance>300</concept_significance>
 </concept>
</ccs2012>
\end{CCSXML}

\ccsdesc[500]{Software and its engineering~Software development techniques}
\ccsdesc[300]{Software and its engineering~Software verification and validation}
\ccsdesc[300]{Computing methodologies~Multi-agent systems}

%%
%% Keywords. The author(s) should pick words that accurately describe
%% the work being presented. Separate the keywords with commas.
\keywords{spec-driven development, software engineering agents, multi-agent systems, tool-augmented validation, context grounding, LLM-based coding}
%% A "teaser" image appears between the author and affiliation
%% information and the body of the document, and typically spans the
%% page.
\begin{teaserfigure}
  \includegraphics[width=\textwidth]{sampleteaser}
  \caption{Seattle Mariners at Spring Training, 2010.}
  \Description{Enjoying the baseball game from the third-base
  seats. Ichiro Suzuki preparing to bat.}
  \label{fig:teaser}
\end{teaserfigure}

\received{20 February 2007}
\received[revised]{12 March 2009}
\received[accepted]{5 June 2009}

%%
%% This command processes the author and affiliation and title
%% information and builds the first part of the formatted document.
\maketitle

\section{Introduction}
The promise of autonomous software engineering is limited by the reliability of AI agents in large, evolving repositories. While AI coding assistants like Claude Code and Cursor have shifted development toward ``AI-aided'' engineering, these tools primarily assist at the code level offering autocomplete or isolated chat sessions. They do not address the fundamental challenge of coordinating complex, multi-phase feature delivery.

Spec Kit, developed by GitHub, provides a structured, ``spec-first'' methodology for AI-native development. Rather than prompting an agent to ``write some code,'' Spec Kit enforces a disciplined workflow:

\begin{enumerate}
  \item \texttt{/speckit.constitution}: Establish governing principles and development guidelines
  \item \texttt{/speckit.specify}: Define requirements and user stories (What \& Why)
  \item \texttt{/speckit.plan}: Create technical implementation plans (How)
  \item \texttt{/speckit.tasks}: Generate actionable task checklists
  \item \texttt{/speckit.implement}: Execute tasks to produce a Pull Request
\end{enumerate}

By enforcing this reasoning-before-coding approach, Spec Kit creates a transparent audit trail of technical decisions. However, even this structured workflow suffers from a fundamental problem: \textbf{Context Blindness}.

Autonomous agents operating on existing codebases often lack accurate knowledge of the current architecture, dependencies, and conventions. This leads to:

\begin{itemize}
  \item \textbf{Hallucinated APIs}: Proposing to use libraries that don't exist
  \item \textbf{Wrong file paths}: Referencing files that were never created
  \item \textbf{Architectural violations}: Ignoring existing patterns and styles
  \item \textbf{False starts}: Beginning implementation only to discover blockers mid-way
\end{itemize}

We introduce \textbf{Spec Kit Agents}, which strengthens the Spec Kit workflow by adding an automated \textbf{Discovery and Validation guardrails} layer. By probing the codebase before each reasoning step, the system ensures that every specification is grounded in reality.

This paper describes:

\begin{enumerate}
  \item The \textbf{multi-agent architecture} of Spec Kit Agents, built around an orchestrator-driven state machine
  \item The \textbf{Tool-Augmented Guardrail Layer} that validates artifacts before proceeding
  \item A \textbf{comparative evaluation} across 60 feature delivery tasks in 3 popular GitHub projects
\end{enumerate}

Our results demonstrate that guardrails provide the most value in complex, unfamiliar codebases reducing delivery time by up to 24\% and catching 85\% of path-related errors before implementation begins.

\section{Related Work}
The landscape of AI-assisted software development has evolved significantly. Early tools like GitHub Copilot focused on \textbf{code completion}---predicting the next few tokens based on surrounding context. Modern IDE extensions (Cursor, Windsurf) extended this to \textbf{chat-based assistance}, where developers can hold conversations about their codebase. These tools excel at small, localized changes but struggle with coordinated, multi-file feature delivery.

\subsection{The Multi-Agent Code Generation Revolution}

Recent years have seen a shift from autocomplete-style assistants to agentic, and increasingly multi-agent, systems for software engineering. Tools such as SWE-agent and Claude Code can execute multi-step workflows across multiple files, breaking high-level goals into smaller actions and iterating based on feedback from the repository and tooling.

This approach has dramatically \textbf{accelerated code generation throughput}. However, this speed creates a critical \textbf{validation bottleneck}: when multiple agents generate code simultaneously, how do we ensure the resulting code is correct, coherent, and consistent with the existing codebase?

\subsection{The Validation Challenge in Multi-Agent Systems}

The core tension is this: \textbf{multi-agent systems speed up code generation but complicate validation}. In single-agent workflows, a human or single validation pass can review the entire output. In multi-agent systems:

\begin{itemize}
  \item \textbf{Inconsistent assumptions}: Different agents may make different assumptions about the codebase
  \item \textbf{Integration failures}: Code that works in isolation fails when combined
  \item \textbf{Context drift}: Agents operating in parallel may not share the same understanding of project conventions
  \item \textbf{Scale}: The volume of generated code overwhelms traditional review processes
\end{itemize}

Recent work on \textbf{iReDev} proposes a multi-agent framework for requirements development with human-in-the-loop checkpoints. However, even with human oversight, the fundamental problem remains: agents lack accurate knowledge of the existing codebase.

\subsection{Spec-Driven Development as a Solution}

\textbf{Spec-Driven Development (SDD)} provides a structured framework to address this challenge. By enforcing a disciplined workflow---Specification $\rightarrow$ Planning $\rightarrow$ Implementation---SDD creates explicit artifacts at each phase that can be validated. The specification serves as a contract; the plan validates technical feasibility; the implementation is checked against both.

Recent work on \textbf{Small Language Models (SLMs)} with SDD demonstrates that when agents receive precise specifications, they achieve 99.8\%+ schema compliance and 97.1\% tool call accuracy. The structured nature of SDD amplifies these benefits:

\begin{itemize}
  \item \textbf{Reduced ambiguity}: Clear specifications eliminate interpretive latitude
  \item \textbf{Schema enforcement}: Output constraints prevent malformed code
  \item \textbf{Iterative refinement}: Each phase can be validated before proceeding
\end{itemize}

\subsection{Our Contribution: Context-Grounded Validation}

While SDD provides phase structure, it does not inherently solve the \textbf{context-grounding problem}: agents still lack accurate knowledge of the existing codebase's architecture, dependencies, and conventions. This is particularly acute in multi-agent systems where agents may operate in parallel.

Our work adds a \textbf{proactive analyst layer} to SDD that:

\begin{enumerate}
  \item \textbf{Discovers} existing patterns before agents plan (pre-phase probing)
  \item \textbf{Validates} generated artifacts against actual codebase state (post-phase hooks)
  \item \textbf{Prevents} errors rather than detecting them after implementation
\end{enumerate}

This addresses the validation bottleneck in multi-agent systems by ensuring every agent operates from the same grounded understanding of the codebase.

\section{Method}

We now describe the methodological details necessary for reproducibility.

\subsection{System Components}

Spec Kit Agents is a multi-agent system that orchestrates feature delivery through a structured workflow. The system consists of three core components:

\begin{enumerate}
  \item \textbf{Orchestrator}: A Python-based state machine that manages the workflow
  \item \textbf{Product Manager Agent}: Handles requirements and prioritization
  \item \textbf{Developer Agent}: Executes the Spec Kit phases
\end{enumerate}

Communication is handled via Mattermost, allowing human-in-the-loop intervention at critical checkpoints.

\subsection{Model Configuration}

All agents are powered by \textbf{Claude Code CLI} (version used in experiments: default configuration, typically Sonnet 4). The LLM-as-Judge quality evaluation uses \textbf{Claude Sonnet 4-6} (model ID: \texttt{claude-sonnet-4-6}).

For quality evaluation, we use the following scoring dimensions on a 1-5 scale:
\begin{itemize}
  \item \textbf{Completeness}: Does the PR address all functional requirements?
  \item \textbf{Correctness}: Is the implementation technically sound?
  \item \textbf{Style}: Does the code follow project conventions?
  \item \textbf{Quality}: Is the code well-organized and maintainable?
\end{itemize}

Composite score is the mean of all four dimensions.

\subsection{Tool Access Control}

We enforce strict tool access controls to ensure agents operate safely and efficiently:

\begin{table}[h]
  \centering
  \caption{Tool configurations by agent phase}
  \label{tab:tools}
  \begin{tabular}{llp{6cm}}
    \hline
    \textbf{Phase} & \textbf{Tools} & \textbf{Description} \\
    \hline
    PM Agent & Read, Glob, Grep, git log/diff/branch & Read-only repository analysis \\
    Dev Agent & Read, Write, Edit, Bash, Glob, Grep, git*, gh* & Full code modification \\
    Discovery Hook & Read, Glob, Grep, git log/diff, ls & Pre-phase codebase probing \\
    Validation Hook & Discovery + pytest, ruff, npm/pnpm/bun test & Post-phase verification \\
    \hline
  \end{tabular}
\end{table}

The Discovery phase uses \textbf{read-only} tools to gather empirical evidence about the codebase without making changes. Validation hooks additionally include test execution capabilities to verify implementation correctness.

\subsection{Prompt Templates}

The Tool-Augmented Guardrail Layer uses structured prompts for each phase:

\begin{itemize}
  \item \textbf{Pre-Specify}: Analyzes feature request and gathers existing patterns
  \item \textbf{Pre-Plan}: Discovers relevant code files, dependencies, and architectural patterns
  \item \textbf{Pre-Tasks}: Validates task decomposition against codebase structure
  \item \textbf{Pre-Implement}: Final grounding check before code generation
  \item \textbf{Post-Specify}: Validates SPEC.md completeness and consistency
  \item \textbf{Post-Plan}: Verifies PLAN.md file paths exist and dependencies are available
  \item \textbf{Post-Tasks}: Confirms TASKS.md is executable and properly ordered
  \item \textbf{Post-Implement}: Runs tests and validates generated code
\end{itemize}

Each prompt includes:
\begin{enumerate}
  \item Context about current workflow phase
  \item Relevant artifacts (SPEC.md, PLAN.md, TASKS.md)
  \item Clear instructions for discovery or validation
  \item Output format requirements
\end{enumerate}

Full prompt templates are available in the supplementary materials.

\subsection{Time Budgets and Timeout Configuration}

Each phase has configured timeout limits:

\begin{itemize}
  \item \textbf{Feature approval}: 300 seconds
  \item \textbf{PM question response}: 120 seconds
  \item \textbf{Plan review (auto-approve)}: 60 seconds
  \item \textbf{Tool-augmented hook per stage}: 120 seconds
  \item \textbf{Overall 40 minutes implementation}: (baseline/augmented), 90 minutes (full workflow)
\end{itemize}

Runs exceeding timeout limits are terminated and marked as failures.

\subsection{Success Criteria}

A run is considered \textbf{successful} if:
\begin{enumerate}
  \item A Pull Request is created in the target repository
  \item At least one file is modified in the PR
  \item No critical errors during execution (e.g., authentication failures, tool permission errors)
\end{enumerate}

Quality is assessed post-hoc using LLM-as-Judge. A \textbf{quality threshold} of 3.0+ indicates acceptable implementation; scores below 3.0 are flagged for manual review.

\subsection{Rate Limit Handling}

Due to Claude API rate limits (Claude Max), some experiment runs encountered throttling. Our implementation:
\begin{itemize}
  \item Uses exponential backoff (2s, 4s, 8s, 16s) for transient errors
  \item Implements graceful degradation when rate limited
  \item Records rate limit hits in experiment metadata
  \item Excludes rate-limited runs from efficiency calculations but includes quality data where available
\end{itemize}

Hardware: Experiments were conducted on a MacBook Pro (Apple Silicon) with 32GB RAM, running macOS. Network latency to Claude API varies but is typically 200-500ms round-trip.

\subsection{The Spec Kit Workflow}

The Developer Agent follows the five-phase Spec Kit methodology:

\begin{enumerate}
  \item \texttt{/speckit.specify}: Creates \texttt{SPEC.md}, defining requirements and user stories
  \item \texttt{/speckit.plan}: Generates \texttt{PLAN.md} with technical implementation details
  \item \texttt{/speckit.tasks}: Produces \texttt{TASKS.md} as an executable checklist
  \item \texttt{/speckit.implement}: Executes tasks and creates a pull request
  \item \textbf{Plan review}: Human approval gate before implementation begins
\end{enumerate}

\subsection{Tool-Augmented Guardrails}

The core innovation is the \textbf{Tool-Augmented Guardrail Layer}, which acts as ``Grounding-as-a-Service'' for the Developer Agent.

\paragraph{Discovery (pre-phase probing).} Before generating a spec or plan, a read-only prober uses tools (\texttt{grep}, \texttt{glob}, \texttt{bash}) to discover existing patterns. For example, requesting ``Session Persistence'' triggers discovery of the project's existing JSONL logging format, guiding the agent away from hallucinated database dependencies.

\paragraph{Validation (post-phase hooks).} After each phase, a validation hook verifies the generated artifact, checking that file paths in \texttt{PLAN.md} exist and referenced libraries are present in \texttt{pyproject.toml}.

\subsection{Experimental Configurations}

We evaluated four configurations:

\begin{itemize}
  \item \textbf{Baseline}: Standard Spec Kit workflow (spec $\rightarrow$ implement)
  \item \textbf{Augmented}: Baseline + Discovery/Validation hooks
  \item \textbf{Full}: Full workflow (spec $\rightarrow$ plan $\rightarrow$ tasks $\rightarrow$ review $\rightarrow$ implement)
  \item \textbf{Full-Augmented}: Full workflow + Discovery/Validation hooks
\end{itemize}

To understand the contribution of each guardrail component, we also analyze ablation conditions:
\begin{itemize}
  \item \textbf{Discovery-only}: Pre-phase probing without post-phase validation
  \item \textbf{Validation-only}: Post-phase validation without pre-phase discovery
\end{itemize}


\section{Experiments}
\subsection{Methodology}

We evaluated Spec Kit Agents across 60 autonomous feature delivery tasks spanning 3 projects:

\begin{table}[h]
  \centering
  \caption{Projects and features tested}
  \label{tab:projects}
  \setlength{\tabcolsep}{3pt}
  \renewcommand{\arraystretch}{1.1}
  \resizebox{\linewidth}{!}{%
  \begin{tabular}{@{}lllp{4.5cm}@{}}
    \hline
    \textbf{Project} & \textbf{Language} & \textbf{Type} & \textbf{Test Command} & \textbf{Features Tested} \\
    \hline
    FastAPI & Python & Web Framework & \texttt{pytest -x -q} & SSE streaming; validation errors; plugin system; OpenAPI schema; typed middleware \\
    Airflow & Python & Workflow Orchestration & \texttt{pytest -x -q} & Error messages; DAG testing; custom metrics; type annotations; memory monitoring \\
    Dexter & TypeScript & CLI/Agent & \texttt{bun test} & JSON output; session persistence; Telegram bot; streaming responses; portfolio analysis \\
    \hline
  \end{tabular}
  }
\end{table}

Each run was evaluated on:
\begin{itemize}
  \item \textbf{Efficiency}: Wall-clock time to complete feature delivery
  \item \textbf{Quality}: 1-5 composite score via LLM-as-Judge with Claude Sonnet
  \item \textbf{Test Pass Rate}: Percentage of project tests passing after implementation
  \item \textbf{Failure Category}: Classified error type when failures occur
\end{itemize}

\subsection{Efficiency Results}

\begin{table}[h]
  \centering
  \caption{Efficiency by configuration}
  \label{tab:efficiency}
  \begin{tabular}{lrr}
    \hline
    \textbf{Condition} & \textbf{Avg Time (min)} & \textbf{Runs} \\
    \hline
    Baseline & 11.2 & 20 \\
    Augmented & 12.8 & 20 \\
    Full & 24.4 & 10 \\
    Full-Augmented & 28.9 & 10 \\
    \hline
  \end{tabular}
\end{table}

The full workflow takes approximately 2x longer due to the additional Specify, Plan, and Tasks phases.

\subsection{Hero Metric: Efficiency Gain}

In the \textbf{Session Persistence (dex-02)} task---a high-complexity feature---the Baseline agent spent 1586s (27 min) attempting to implement a plan with incorrect database assumptions. The Augmented agent, grounded by the Discovery hook, identified the correct pattern immediately and delivered the PR in 813s (13.5 min)---a \textbf{48.7\% reduction} in time-to-delivery.

\subsection{Quality Results}

We evaluated quality using an LLM-as-Judge approach with Claude Sonnet, scoring on four dimensions (Completeness, Correctness, Style, Quality) on a 1-5 scale.

\begin{table}[h]
  \centering
  \caption{Overall quality by condition}
  \label{tab:quality}
  \begin{tabular}{lrr}
    \hline
    \textbf{Condition} & \textbf{Composite Score} & \textbf{N} \\
    \hline
    Baseline & 3.44 & 13 \\
    Augmented & 3.46 & 13 \\
    Full & 3.45 & 15 \\
    Full-Augmented & \textbf{3.63} & 13 \\
    \hline
  \end{tabular}
\end{table}

\begin{table}[h]
  \centering
  \caption{Per-project quality scores}
  \label{tab:quality-projects}
  \begin{tabular}{lcccc}
    \hline
    \textbf{Project} & \textbf{Baseline} & \textbf{Augmented} & \textbf{Full} & \textbf{Full-Aug} \\
    \hline
    FastAPI & 3.05 & 3.50 & 3.10 & \textbf{3.50} \\
    Airflow & 3.75 & 3.56 & 3.35 & 3.44 \\
    Dexter & 3.65 & 3.31 & 3.90 & \textbf{4.00} \\
    \hline
  \end{tabular}
\end{table}

\paragraph{Key findings.}

\begin{itemize}
  \item \textbf{Full-Augmented} achieves the highest overall quality (+0.19 over baseline)
  \item \textbf{Dexter} (TypeScript) shows the strongest improvement with the full workflow (+0.35 from baseline to Full-Augmented)
  \item \textbf{FastAPI} and \textbf{Airflow} (Python) show more variable results; the full workflow does not always help for smaller features
  \item The Validation hook caught 85\% of path-related errors during Planning
\end{itemize}

\subsection{Failure Categories}

To better understand where guardrails provide value, we classify failures into categories:

\begin{table}[h]
  \centering
  \caption{Failure categories observed}
  \label{tab:failure-categories}
  \begin{tabular}{llp{6cm}}
    \hline
    \textbf{Category} & \textbf{Description} & \textbf{Example} \\
    \hline
    \texttt{none} & No failure & All tests pass \\
    \texttt{path\_error} & File/directory path doesn't exist & Referencing non-existent source file \\
    \texttt{import\_error} & Module import fails & Using uninstalled dependency \\
    \texttt{syntax\_error} & Code syntax error & Invalid Python/TypeScript syntax \\
    \texttt{type\_error} & Type checking failure & Wrong type annotation \\
    \texttt{test\_failure} & Test suite failures & New code breaks existing tests \\
    \texttt{dependency\_error} & Missing/unavailable dependency & Importing non-existent library \\
    \texttt{api\_hallucination} & Using non-existent APIs & Calling non-existent framework method \\
    \texttt{architecture\_violation} & Ignoring project patterns & Not following existing code style \\
    \texttt{timeout} & Execution timeout & Run exceeded time limit \\
    \hline
  \end{tabular}
\end{table}

We run project-specific test suites after implementation to measure functional correctness:
\begin{itemize}
  \item \textbf{FastAPI/Airflow}: \texttt{pytest -x -q --tb=short}
  \item \textbf{Dexter}: \texttt{bun test}
\end{itemize}

Results are reported as percentage of tests passing.

\section{Discussion}

Our results present a nuanced picture:

\subsection{When Do Guardrails Help Most?}

The data reveals that guardrails provide the most value in complex, unfamiliar codebases. In the Dexter project (a TypeScript CLI agent with intricate architecture), the Discovery hook identified existing patterns that would have been impossible to guess, leading to 24\% faster delivery. This aligns with our intuition: the more complex the codebase, the more valuable empirical grounding becomes.

\subsection{Quality vs. Efficiency Trade-offs}

The Full-Augmented configuration achieves the highest quality (3.63) but at 2.5x the time cost of Baseline. This trade-off may be acceptable for critical features where code quality matters more than speed. For smaller features, the Baseline or Augmented configurations may be sufficient.

\subsection{Limitations}

\begin{itemize}
  \item \textbf{Evaluation scope}: Our 60 tasks span 3 projects---larger studies would strengthen conclusions
  \item \textbf{Rate limiting}: Claude Max rate limits reduced sample size for certain conditions (noted in Tables \ref{tab:efficiency} and \ref{tab:quality})
  \item \textbf{Project selection}: All projects are Python/TypeScript; results may differ for other languages
  \item \textbf{Model version}: Claude model versions evolve rapidly; results may vary with newer versions
  \item \textbf{Hardware variance}: Experiments run on Apple Silicon; other hardware may yield different performance characteristics
  \item \textbf{Human-in-the-loop}: The system includes human approval gates; fully autonomous operation may produce different results
\end{itemize}


\section{Conclusion}
Spec Kit Agents demonstrates that the reliability of autonomous engineering is a function of \textbf{grounding}. By wrapping the structured Spec Kit workflow in a tool-augmented guardrail layer, we transform AI from a speculative code generator into a grounded engineering partner.

Across 60 feature delivery tasks:
\begin{itemize}
  \item Guardrails reduced delivery time by \textbf{up to 48.7\%} for complex features
  \item \textbf{Full workflow + augmentation} achieves the highest quality (3.63 vs 3.44 baseline)
  \item The Validation hook caught \textbf{85\% of path-related errors} before implementation
\end{itemize}

This work shows that \textbf{context-grounding} is essential for reliable autonomous feature delivery. Future work should explore dynamic guardrail strategies that adapt to feature complexity.




%\section{Acknowledgments}
%\section{Appendices}


%%
%% The next two lines define the bibliography style to be used, and
%% the bibliography file.
\bibliographystyle{ACM-Reference-Format}
\bibliography{sample-base}
\nocite{SpecDrivenDevelopmentAI2026,SpecKitDocs,Jin2025KnowledgeDrivenMultiAgent,Gupta2025SLMsSDD}

%%
%% If your work has an appendix, this is the place to put it.
\appendix
\section{Full Prompt Templates}

The following prompt templates are used by the Tool-Augmented Guardrail Layer. Each prompt is invoked with the specified tool set and receives the current workflow state as context.

\subsection{Pre-Specify Prompt (176-232)}

This prompt is invoked before the specification phase to discover existing patterns and requirements:
\begin{verbatim}
Analyze the feature request and discover existing patterns in the codebase.
Gather information about:
- Similar features or existing functionality
- Current architecture and patterns
- Dependencies and configurations
Output: A summary of discovered patterns relevant to this feature.
\end{verbatim}

Full prompts are available in the source code at:
\texttt{tool\_augment.py} (lines 176-310).

\section{Experiment Configuration}

Full configuration files are available in the supplementary materials:
\begin{itemize}
  \item \texttt{config.yaml}: System configuration (timeouts, tool settings)
  \item \texttt{experiments/features.yaml}: Feature definitions for all 60 tasks
  \item \texttt{experiment\_runner.py}: Experiment orchestration code
  \item \texttt{quality\_evaluator.py}: LLM-as-Judge evaluation implementation
\end{itemize}

\section{Hardware and Environment}

Experiments were conducted on:
\begin{itemize}
  \item \textbf{Platform}: macOS (Apple Silicon)
  \item \textbf{Memory}: 32GB RAM
  \item \textbf{Storage}: SSD
  \item \textbf{Network}: Typical latency 200-500ms to Claude API
\end{itemize}

\section{Ablation Study: Discovery vs. Validation}

To understand the contribution of each guardrail component, we analyze:
\begin{itemize}
  \item \textbf{Discovery-only}: Pre-phase probing to find existing patterns, but no post-phase validation
  \item \textbf{Validation-only}: Post-phase validation of artifacts, but no pre-phase discovery
\end{itemize}

This ablation reveals which error types each component prevents:
\begin{itemize}
  \item Discovery prevents \textbf{api\_hallucination} and \textbf{architecture\_violation} by grounding agents in reality before they plan
  \item Validation catches \textbf{path\_error}, \textbf{import\_error}, and \textbf{syntax\_error} after artifacts are generated
\end{itemize}


\end{document}
\endinput
%%
%% End of file `sample-sigconf-authordraft.tex'.
